\section{Método}

\subsection{Diseño}
Investigación aplicada, enfoque cuantitativo, nivel explicativo-comparativo.
Diseño experimental comparativo para la fase técnica y cuasiexperimento de medidas
repetidas para la fase pedagógica.

\subsection{Corpus y muestra}
Se conforma un corpus de 40 páginas manuscritas en Braille Grado 1 (35
controladas y 5 espontáneas), capturadas bajo 12 combinaciones fotométricas, lo
que produce 480 fotografías. La partición es por página y se aplica validación
cruzada de cinco pliegues agrupada. La fase pedagógica prevé hasta 30 docentes de
las UGEL Ilo y Mariscal Nieto.

\subsection{Variables e instrumentos}
La variable independiente es la configuración del detector (YOLOv8n o YOLOv11;
FP32 o INT8). La variable dependiente es el desempeño: precisión (mAP@0.5, F1
macro, Kappa, AUC-ROC, CER, WER) y eficiencia (latencia, memoria, tamaño). La
oportunidad de la retroalimentación se mide con el tiempo de revisión por página;
la carga cognitiva y la usabilidad con NASA-TLX \parencite{hart1988} y SUS
\parencite{brooke1996}.

\subsection{Procedimiento y análisis}
Se aplican Wilcoxon, McNemar, bootstrap de 5000 remuestreos y un criterio de no
inferioridad para H2. Se documentan versiones de software, hardware y semillas
para asegurar la reproducibilidad.
